\section{Descripción del Problema}

En mercados de criptoactivos, la manipulación coordinada mediante esquemas de
\emph{pump \& dump} representa un desafío significativo para la detección
temprana de anomalías. En estos eventos, un grupo de actores acumula un activo
de baja liquidez y luego induce de forma coordinada un aumento artificial de su
precio para finalmente vender en el pico. Este proceso ocurre en escalas de
tiempo cortas y suele dejar señales débiles antes de la fase de \emph{pump}, lo
que dificulta su identificación anticipada.

A diferencia de otros fenómenos financieros, estos eventos no se manifiestan
únicamente en el comportamiento individual de un activo, sino en patrones
colectivos. Múltiples tokens pueden exhibir movimientos sincronizados como
resultado de estrategias coordinadas de acumulación, rotación de capital o
manipulación cruzada. El problema central consiste en detectar, a partir de
series temporales de precios, la existencia de un subconjunto desconocido de
tokens cuyos retornos presentan correlación estadística atípica en una ventana
temporal previa a un evento de \emph{pump} ---subconjunto no etiquetado a
priori, de tamaño variable, cuyo impacto puede quedar oculto dentro del
comportamiento agregado del mercado.

\subsection{Contexto del mercado}

Los mercados de criptoactivos funcionan las 24 horas, tienen miles de tokens
activos y poca supervisión. La mayoría de estos tokens tienen una capitalización
muy baja, lo que significa que un grupo con capital moderado puede mover su
precio fácilmente. Los organizadores de un \emph{pump \& dump} coordinan la
compra a través de canales de Telegram o Discord, inflan el precio y venden
antes de que los demás puedan reaccionar~\cite{lamorgia2023}.

Este problema presenta tres dificultades fundamentales. Primero, las
correlaciones estimadas en ventanas finitas están fuertemente contaminadas
por ruido estadístico, lo que dificulta distinguir entre correlación genuina
y fluctuaciones aleatorias. Segundo, el mercado de criptoactivos es altamente
dinámico: las relaciones entre activos cambian rápidamente debido a factores
externos (noticias, liquidez, choques de mercado), lo que complica la
definición de una línea base estable. Tercero, los eventos de \emph{pump \&
dump} ocurren en escalas temporales cortas ---horas o pocos días---, lo que
limita la cantidad de datos disponibles para estimar de forma robusta la
estructura de correlación.

Un problema técnico importante es que en condiciones normales, casi todos los
tokens se mueven juntos siguiendo a Bitcoin ---lo que se llama el \emph{modo de
mercado global}~\cite{plerou1999}. Si no eliminamos esta correlación de fondo,
la red de correlaciones queda tan densa que no se puede distinguir ninguna
estructura: todo parece conectado con todo.

\subsection{Pregunta central}

Antes de que un \emph{pump} sea visible en el precio, los organizadores acumulan
posiciones de forma coordinada. Esta acumulación debería dejar una huella
estructural en la red de correlaciones entre tokens: un subgrupo que de repente
se mueve junto, de manera distinta al resto del mercado.

La pregunta que guía este trabajo es:
\begin{quote}
  \emph{¿Es posible detectar la fase de acumulación previa a un \emph{pump \&
  dump} observando cambios en la distribución de localización espectral de la
  red de correlaciones entre tokens?}
\end{quote}

Responder esto requiere construir un modelo que formalice qué significa
``localización espectral'' y cómo medirla de forma robusta en redes dinámicas.
Ese es el objetivo de las secciones~\ref{sec:formalizacion} y~\ref{sec:modelo}.

\subsection{Limitaciones de los enfoques actuales}

Los métodos que analizan precio y volumen de cada token por separado tienen dos
problemas principales. Primero, generan muchas falsas alarmas: cualquier
movimiento brusco legítimo puede parecerse a un \emph{pump}. Segundo, no ven la
coordinación entre tokens, que es justo donde está la huella de la
manipulación.

Los métodos basados en grafos que existen se enfocan en otras cosas: detectar
ciclos de arbitraje~\cite{wang2022} o analizar flujos entre
exchanges~\cite{soramaki2007}. Si bien el IPR se ha calculado sobre matrices de
correlación financiera antes~\cite{plerou1999,nobi2013}, ningún trabajo lo usa
como detector de manipulación combinando: (i)~filtrado de ruido con RMT antes
de construir el grafo; (ii)~estadísticos sobre todo el espectro, sin asumir en
qué región aparece la anomalía; y (iii)~validación sobre eventos confirmados
con hora exacta. Esa es la contribución de este trabajo.
